\documentclass[10.5pt,a4paper]{article}
\usepackage[a4paper,top=1.3cm,bottom=1.35cm,left=1.35cm,right=1.35cm]{geometry}
\usepackage{fontspec}
\usepackage{xeCJK}
\usepackage{enumitem}
\usepackage[hidelinks]{hyperref}
\usepackage{tabularx}
\usepackage{array}
\usepackage{xcolor}
\usepackage{titlesec}
\usepackage{fancyhdr}
\usepackage{lastpage}

% 页面与页脚设置（针对纸质打印优化双面/单面页码标识）
\pagestyle{fancy}
\fancyhf{}
\renewcommand{\headrulewidth}{0pt}
\cfoot{\footnotesize\color{subtextgray}第 \thepage{} 页 \ / \ 共 \pageref{LastPage} 页}

% 字体配置（与现有环境保持一致）
\setmainfont{texgyrepagella-regular.otf}[
  BoldFont = texgyrepagella-bold.otf,
  ItalicFont = texgyrepagella-italic.otf,
  BoldItalicFont = texgyrepagella-bolditalic.otf
]
\setCJKmainfont{FandolSong-Regular.otf}[
  BoldFont = FandolSong-Bold.otf,
  ItalicFont = FandolKai-Regular.otf
]
\setCJKsansfont{FandolHei-Regular.otf}[
  BoldFont = FandolHei-Bold.otf
]
\setCJKmonofont{FandolFang-Regular.otf}

% 打印专用色彩配置（高对比度，黑白及灰阶打印友好）
\definecolor{titlecolor}{HTML}{142634}
\definecolor{textmain}{HTML}{111111}
\definecolor{textgray}{HTML}{333333}
\definecolor{subtextgray}{HTML}{555555}
\definecolor{rulecolor}{HTML}{CCCCCC}

% 标题格式与间距控制
\titleformat{\section}{\large\bfseries\color{titlecolor}}{}{0pt}{}[\vspace{1.5pt}{\color{rulecolor}\hrule height 0.6pt}]
\titlespacing*{\section}{0pt}{6pt}{3.5pt}

% 列表与段落排版微调
\setlist[itemize]{leftmargin=1.35em, itemsep=1.5pt, topsep=1.5pt, parsep=0pt}
\setlength{\parindent}{0pt}
\setlength{\parskip}{2pt}

% 自定义宏命令
\newcommand{\entry}[4]{%
  \noindent\textbf{#1} \hfill #2\\
  {\color{subtextgray}#3} \hfill {#4}\par\vspace{1.5pt}
}
\newcommand{\textlink}[2]{\href{#1}{\color{textmain}#2}}

\begin{document}
\color{textmain}

% ==================== 个人头部信息 ====================
{\LARGE\bfseries\color{titlecolor} 牛群}\\[3pt]
北京理工大学 \quad 机械工程博士研究生（预计 2026 年 12 月毕业） \quad 目标岗位：企业研发 / 博士后\\
邮箱：\textlink{mailto:nq_eecm@163.com}{nq\_eecm@163.com} \quad 手机：15601280520 \quad 个人网站：\textlink{https://beta1scat.github.io/}{beta1scat.github.io} \quad 学术主页：\textlink{https://scholar.google.com/citations?hl=en&user=6E4Hd8UAAAAJ}{Google Scholar}

% ==================== 个人简介 ====================
\section*{个人简介}
在北京理工大学攻读机械工程博士学位期间，聚焦复杂非结构化场景中的机器人感知、几何建模与精细操作决策。系统开展三维视觉、解析运动学、高精度手眼标定、点云位姿估计及语言引导抓取等算法链路研究，致力于消除误差累积并打通“感知--建模--规划--执行”的闭环系统；在夯实几何先验与动力学解析方法的同时，深入探索具身智能与大模型赋能下的机器人自主技能学习。工程实践涵盖七自由度双臂协同系统、工业无序分拣、自适应浮动曲面打磨及仿生垂直越障装备等研发，精通 C/C++、Python、PyTorch 与 ROS/ROS 2。

% ==================== 教育经历 ====================
\section*{教育经历}
\entry{北京理工大学}{2020.09 -- 至今}{机械工程}{博士（在读）}
\entry{北京工业大学}{2017.09 -- 2020.06}{机械工程}{硕士}
\entry{河北科技大学}{2013.09 -- 2017.06}{过程装备与控制工程}{学士}

% ==================== 研究方向 ====================
\section*{研究方向}
\begin{itemize}
	\item \textbf{机器人感知与场景理解}：杂乱场景实例分割、点云语义分类、几何基元自适应拟合及工业 3D 视觉系统集成。
	\item \textbf{机构运动学与标定}：同构机器人通用解析逆运动学建模、手眼关系与工具坐标系精密联合标定。
	\item \textbf{视觉引导抓取}：点云位姿估计算法超参数自适应优化、物理稳定性约束下的抓取位姿生成与几何干涉过滤。
	\item \textbf{具身操作与语言引导}：结合多模态空间定位、三维几何推理与任务时序约束的开放场景机器人抓取规划。
\end{itemize}

% ==================== 代表性学术论文 ====================
\section*{代表性学术论文}
\begin{itemize}
	\item \textbf{Niu, Qun}, Zhang, Chuanlin, Zhang, Tianyu, Zhao, Jieliang, Fu, Tie, and Chen, Xuemei. Language-Guided Robot Grasping Based on Basic Geometric Shape Fitting. \textit{Advanced Intelligent Systems}, e202501276, 2025.\\
	\textit{\small 贡献：提出融合多模态目标定位与基本几何体拟合的语言引导抓取框架，显著提升无模型条件下的抓取准确率。}
	\item \textbf{Niu, Qun}, Zhao, Jieliang, Liang, Lulu, Xing, Jin, Li, Hongkun, and Wang, Ziru. A General Framework for the Analytical Inverse Kinematics Solution of Industrial Robots. \textit{Proc. IMechE, Part C: J. Mech. Eng. Sci.}, 239(12):4499--4511, 2025.\\
	\textit{\small 贡献：构建同构工业机器人通用的解析逆运动学求解框架，使异构同类构型共享标准化 D-H 闭式解法。}
	\item \textbf{Niu, Qun}, Wang, Ziru, Li, Hongkun, and Zhao, Jieliang. A Parameters Optimization Framework for Pose Estimation Algorithm Based on Point Cloud. \textit{Journal of Physics: Conference Series}, 2746(1):012039, 2024.\\
	\textit{\small 贡献：建立面向点云位姿估计的超参数自动化优化机制，大幅降低不同几何工件落地中的人工调优成本。}
	\item \textbf{Niu, Qun}, Zhao, Jieliang, Liang, Lulu, and Xing, Jin. Dynamic Vertical Climbing Mechanism of Chinese Dragon-Li Cats Based on the Linear Inverted Pendulum Model. \textit{Journal of Bionic Engineering}, 20(1):136--145, 2023.\\
	\textit{\small 贡献：基于倒立摆模型解析家猫动态垂直跃升机理，揭示非静态接触条件下的仿生缓冲与高效附着规律。}
	\item \textbf{Niu, Qun}, et al. A Simultaneous Hand-Eye and Tool Coordinate System Calibration Framework for Robotic Grasping. \textit{Intelligent Service Robotics} (under review).\\
	\textit{\small 贡献：建立手眼几何关系与工具坐标系一体化同步标定模型，有效切断分步标定引发的误差传递链条。}
	\item \textbf{Niu, Qun}, et al. Enhancing Scene Perception: Shape Fitting with Geometric Primitives. \textit{Intelligent Service Robotics} (under review).\\
	\textit{\small 贡献：融合 SAM 模型与三维几何基元分解技术，实现非结构化杂乱场景的高效结构化几何感知。}
\end{itemize}

% ==================== 发明专利 ====================
\section*{发明专利}
\begin{itemize}
	\item 牛群，赵杰亮等. 基于视觉引导的机器人系统及机器人系统校准方法：授权，CN202410192128.6[P].
	\item 牛群，赵杰亮等. 用于机器人的关节角度生成方法及机器人系统：授权，CN202310966039.8[P].
	\item 赵杰亮，牛群等. 一种高度和角度可调节的双滑块式斜坡实验台：授权，CN202110431126.4[P].
	\item 牛群，赵杰亮等. 位姿估计模型配置方法及位姿估计方法：实审，CN202311698845.8[P].
	\item 其他成果：四足仿生与爬壁机器人合作学术论文 9 篇；另获机器人控制与机构相关授权或实审专利 6 项。
\end{itemize}

% 严格强制换页：确保第 1 页为学术背景与学术产出，第 2 页为项目落地、工程实践与技能荣誉
\newpage

% ==================== 在校科研与重点工程项目 ====================
\section*{在校科研与工程项目}
\entry{面向抓取任务的目标检测与路径规划算法研究}{2022.08 -- 2024.08}{项目骨干}{在校项目}
主导无序分拣场景全流程抓取链路设计与落地：实现机械臂多自由度运动学求解、笛卡尔空间路径规划及实时碰撞干涉规避；开发高鲁棒性手眼标定与点云位姿估计模型；设计融合几何包络特征与物理接触稳定性的抓取位姿快速生成与多指标过滤算法。

\entry{基于一体化关节的仿生四足越障机器人研究}{2020.09 -- 2022.08}{项目骨干}{在校项目}
负责基于昆虫腿部结构的一体化高爆发关节驱动器研发与仿生动态越障动力学建模；搭建高速摄像实验系统采集活体动态跃升姿态序列，基于 SLIP 动力学模型深入解析机体冲击吸收与动态跃迁机理，完成具备主被动弹性储能特征的驱动样机设计。

\entry{双臂机器人拟人化动作协调控制方法研究}{2017.09 -- 2020.06}{项目骨干}{在校项目}
针对线驱动七自由度轻量化双臂机器人，建立高精度 D-H 机构学模型；提出面向拟人化协调操作的逆运动学闭式求解算法，并针对精密轴孔装配工艺规划了平滑的关节空间与任务空间协调控制轨迹。

% ==================== 产业与实习经历 ====================
\section*{产业与实习经历}
\entry{北京银河通用机器人有限公司}{2025.08 -- 2026.02}{算法工程实习生}{具身智能研发}
深入参与银河通用太空舱项目，依托轮式双臂具身机器人 G1，研发商业零售场景下的自适应抓取与移动操作技能；系统结合解析几何抓取先验与现代学习范式，实现机器人在复杂多变货架及狭窄作业空间中的自主泛化抓取与安全作业。

\entry{北京迁移科技有限公司}{2022.08 -- 2024.08}{研发部算法实习生}{工业 3D 视觉与机器人}
负责工业机械臂运动学求解内核、轨迹优化及碰撞检测模块的研发维护；主导抓取位姿自适应生成、非注册通用抓取算法预研与手眼自动标定工具链；基于 Blazor 框架开发 3D 场景配置与几何数字孪生交互系统，提升产线快速部署效率。

\entry{中国科学院自动化研究所}{2019.04 -- 2020.05}{数字内容技术与服务研究中心}{机器人系统研发}
研发打磨机器人复杂二次曲面自适应轨迹规划算法及力控浮动磨头闭环控制程序，实现工业现场精密交付；开发多轴同步控制与基于 PID 控制器的螺栓精准预紧力装配系统；完成模块化防水蛇形机器人机构与电控布局设计。

\entry{北京墨狄科技有限公司}{2018.04 -- 2019.04}{研发部实习生}{机械臂与机电控制}
参与四自由度仿生连杆机械臂正逆运动学建模、奇异点规避及关节空间轨迹规划；参与多工作模式细胞培养生物反应器机电系统研发，实现上位机交互系统与高精度磁力搅拌电控驱动的协调控制。

% ==================== 专业技能 ====================
\section*{专业技能}
\begin{itemize}
	\item \textbf{编程与工具栈}：C/C++、Python、PyTorch、ROS / ROS 2、Linux、Git、Blazor、MATLAB。
	\item \textbf{机器人控制与运动学}：正逆运动学解析建模、轨迹规划与平滑插补、碰撞检测、手眼与工具坐标系标定、抓取位姿生成与过滤。
	\item \textbf{机器视觉与空间感知}：3D 视觉系统集成、点云处理（PCL/Open3D）、6D 位姿估计、实例分割、空间几何基元拟合。
	\item \textbf{前沿学习范式与具身智能}：视觉-语言-动作模型（VLA）、强化学习（RL）在机械臂操作中的应用实践与前沿跟踪。
\end{itemize}

% ==================== 荣誉与获奖 ====================
\section*{荣誉与获奖}
\begin{itemize}
	\item 北京工业大学优秀研究生奖、研究生学习优秀一等奖、研究生科研优秀奖。
	\item 第八届河北省大学生机械设计创新大赛二等奖。
	\item 体育竞技：篮球赛事市级奖项 1 项、校级奖项 5 项。
\end{itemize}

\end{document}
